

\begin{abstract}
仅凭 F3、Fz、F4 三路额区脑电，既要从伪迹中提取视觉响应，又要解释左右三角形刺激产生的差异，困难在于有效信号较弱、脑源与电极之间的对应关系不唯一。本文利用两组视觉认知实验的四份记录，依次建立试次质量筛选、视觉刺激正向计算和认知功能代理模型，分析三通道观测能够支持的波形特征与机制解释。

\textbf{针对问题一，}采用稳健黎曼伪迹拒绝方法，结合协方差距离和分频段统计量构造信号质量指数 SQI。位置提示任务配置四个低维检测器，形状提示任务配置三个检测器并放宽筛选阈值；剔除削顶及低质量试次后，进行基线校正、条件叠加平均和候选峰提取。再以二阶差分惩罚的三次 B 样条拟合 ERP，由广义交叉验证选择平滑参数。四份记录共保留 \textbf{297 个试次}，24 条条件平均曲线的拟合优度 \(R^2\) 为 \(0.891\)--\(1.000\)。阈值敏感性分析与等样本量随机抽样表明，整体波形相对稳定，但局部峰值仍受筛选强度影响。

\textbf{针对问题二，}将 LGN 的 ON/OFF 中继、Gabor 方向滤波、三角形空间构型编码与 Wilson--Cowan 型兴奋--抑制群体方程连接，构造五个功能源代理，并通过固定线性映射生成三通道脑电。采用有界 L-BFGS-B 搜索时间常数与抑制增益，左右条件共享幅度因子。空间位置消融后，模型左右差异的均方根幅度降至原来的 \(10^{-5}\) 量级，说明位置编码是该模型产生差异的关键。总体 ERP 的样本内相关系数为 \textbf{0.894}，归一化均方根误差为 \(0.358\)；左右差异波的相关系数仅为 \(-0.097\)。基于实测脑电九维窗口特征的收缩线性判别分析，在留一记录验证中取得 \(46.6\%\) 的平衡准确率，尚未获得稳定的跨记录判别能力。

\textbf{针对问题三，}直接从原始脑电提取 ERP 均值与频带对数功率，构造视觉前馈、记忆匹配和认知控制三类功能代理 \(V/H/P\)，以普通最小二乘描述试次内关联，并通过留一记录验证、代理消融及 Logistic 应答侧预测检验模型。按提示后 \(2.2\,\mathrm{s}\) 的候选目标时刻对齐后，加入 \(H/P\) 代理及任务变量，特征重建的标准化 RMSE 由 \textbf{1.015 降至 0.856}；13 维脑电特征的提示方向分类平衡准确率为 \(52.2\%\)。这些结果支持代理对观测特征的压缩表达，但不足以确认深部脑源或认知通路的因果关系。

\keywords{脑电图\quad 稳健黎曼伪迹拒绝\quad 惩罚样条\quad 神经群体动力学\quad 认知功能代理}
\end{abstract}

